% Part: history 
% Chapter: biographies 
% Section: julia-robinson
\documentclass[../../../include/open-logic-section]{subfiles}

\begin{document}

\olfileid{his}{bio}{rob}

\olsection{جولیا رابنسن}

\olphoto{robinson-julia}{جولیا رابنسن}

جولیا بومین رابنسن امریکی ریاضی دان سی۔ اوہ خاص کر کے فیصلہ کرن والے مسئلیاں اُتے اپنے کم لئی جانی جاندی اے، تے سبھ توں ودھ ہلبرٹ دے دسویں مسئلے دے حل وچ اپنے حصے لئی۔ رابنسن 8 دسمبر 1919 نوں سینٹ لوئس، میزوری وچ جمی۔ رابنسن نوں چیتا سی کہ اوہ بچپن توں ای عدداں وچ دلچسپی لیندی سی \citep[4]{Reid1986}۔ نو سال دی عمر وچ اوہنوں لال بخار ہویا تے فیر گٹھیا دے بخار دے کئی دورے پئے۔ ایس کرکے اوہنوں بہت ویلا منجے اُتے گزارنا پیا تے اوہ پڑھائی وچ پچھے رہ گئی۔ نجی استاداں دی مدد نال اوہ پڑھائی وچ فیر نال رَل گئی، پر بیماری دے جسمانی اثر اوہدی ساری زندگی اُتے رہے۔

بچپن دیاں اوکڑاں دے باوجود رابنسن نے ریاضی تے سائنس دے کئی انعاماں نال ہائی سکول پورا کیتا۔ یونیورسٹی دی پڑھائی اوہنے سین ڈیاگو سٹیٹ کالج توں شروع کیتی تے آخری سال برکلی دی یونیورسٹی آف کیلیفورنیا چلی گئی۔ اوتھے ریاضی دان رافائل رابنسن دا اوہدے اُتے اثر پیا۔ اوہ چنگے دوست بن گئے تے 1941 وچ ویاہ کر لیا۔ استاداں وچوں اک دی گھر والی ہون کرکے رابنسن نوں برکلی دے ریاضی شعبے وچ پڑھاؤن دی کھُل نہ ملی۔ بھاویں اوہ بناں داخلے ریاضی دیاں جماعتاں سُندی رہی، اوہ یونیورسٹی چھڈ کے ٹبّر شروع کرنا چاہندی سی۔ پر ویاہ توں تھوڑی دیر پچھوں رابنسن نوں نمونیا ہو گیا۔ اوہنوں دسیا گیا کہ بچپن دے گٹھیا بخار کارن اوہدے دل اُتے زخم دے نشان والا ٹشو کافی جم چکیا سی۔ ایس ٹشو دی سنگینی کارن ڈاکٹر نے اندازہ لایا کہ اوہ چالی سال توں اگے نہیں جی سکدی، تے اوہنوں بچے نہ کرن دی صلاح دتی گئی \citep[13]{Reid1986}۔

رابنسن لمّا سماں افسردہ رہی، پر آخر ریاضی دی پڑھائی جاری رکھن دا فیصلہ کیتا۔ اوہ برکلی مُڑ آئی تے 1948 وچ الفریڈ تارسکی دی نگرانی تھلّے پی ایچ ڈی پوری کیتی۔ تارسکی وکھا چکیا سی کہ حقیقی عدداں دے پہلے درجے دے نظریے دا فیصلہ ہو سکدا اے، تے گوڈل دے کم توں نکلدا سی کہ قدرتی عدداں دے پہلے درجے دے نظریے دا فیصلہ نہیں ہو سکدا۔ ایہ اک وڈا کھُلّا سوال سی کہ ناطق عدداں دے پہلے درجے دے نظریے دا فیصلہ ہو سکدا اے یا نہیں۔ اپنے مقالے \citeyearpar{Robinson1949} وچ رابنسن نے ثابت کیتا کہ ایہدا فیصلہ نہیں ہو سکدا۔

فیصلہ کرن والے مسئلیاں وچ دلچسپی کرکے رابنسن نے اگے ہلبرٹ دا دسواں مسئلہ حل کرن دی کوشش کیتی۔ ایہ مسئلہ 23 مشہور ریاضیاتی مسئلیاں دی اوہ فہرست وچوں سی جیہڑی ڈیوڈ ہلبرٹ نے 1900 وچ پیش کیتی سی۔ دسواں مسئلہ پُچھدا اے کہ کی کوئی ایہو جیہا الگورتھم اے جیہڑا محدود سمیں وچ جواب دے سکے کہ صحیح عددی ضریباں والی کثیر رکنی مساوات، جیویں $3x^2 - 2y +3 = 0$، دا صحیح عدداں وچ حل اے یا نہیں۔ ایہو جیہے سوالاں نوں \emph{ڈیوفانتی مسئلے} آکھدے نیں۔ کجھ مُڈھلیاں کامیابیاں پچھوں رابنسن مارٹن ڈیوس تے ہلیری پٹنم نال رَل گئی، جیہڑے ایس مسئلے اُتے کم کردے پئے سن۔ اوہناں نے وکھایا کہ گھات والے ڈیوفانتی مسئلیاں (جتھے انجان عدد گھات دی تھاں وی آ سکدے نیں) دا فیصلہ نہیں ہو سکدا۔ اوہناں نے وی وکھایا کہ اک خاص قیاس (جیہنوں بعد وچ «J.R.» آکھیا گیا) توں لازم آؤندا اے کہ ہلبرٹ دے دسویں مسئلے دا فیصلہ نہیں ہو سکدا \citep{DavisPutnamRobinson1961}۔ رابنسن 1960 دی ساری دہائی ایس مسئلے اُتے کم کردی رہی۔ 1970 وچ نوجوان روسی ریاضی دان یوری ماتیاسےوِچ نے آخر J.R. دا مفروضہ ثابت کر دتا۔ رَلواں نتیجہ ہُن ماتیاسےوِچ–رابنسن–ڈیوس–پٹنم قضیہ، یا مختصر MRDP قضیہ، کہلاندا اے۔ ماتیاسےوِچ تے رابنسن دوست بن گئے تے کئی مقالیاں اُتے رَل کے کم کیتا۔ ماتیاسےوِچ نوں اک خط وچ رابنسن نے اک واری لکھیا، «مینوں سچ مُچ بڑی خوشی اے کہ رَل کے کم کردیاں (ہزاراں میل دور رہندیاں) اسیں صاف ای اوہتوں ودھ اگے ودھ رہے آں جِنّا ساڈے وچوں کوئی اک اکلّا ودھ سکدا سی» \citep[45]{Matijasevich1992}۔

رابنسن امریکن میتھمیٹیکل سوسائٹی دی پہلی عورت صدر سی تے نیشنل اکادمی آف سائنس لئی چُنی جان والی پہلی عورت سی۔ لیوکیمیا دی تشخیص پچھوں اوہ 30 جولائی 1985 نوں 65 سال دی عمر وچ وفات پا گئی۔

\begin{reading}
رابنسن دے ریاضیاتی مقالے اوہدیاں \textit{Collected
  Works} وچ ملدے نیں \citep{Robinson1996}۔ اوہدے وچ نیشنل اکادمی آف سائنسز لئی اوہدی سوانح دا مُڑ چھاپیا روپ وی اے \citep{Feferman1994}۔ رابنسن دی وڈی بھین کونسٹینس ریڈ نے انٹرویوواں دی بنیاد اُتے «Autobiography of Julia» چھپوائی \citep{Reid1986} تے پوری یادداشت وی لکھی \citep{Reid1996}۔ رابنسن تے ہلبرٹ دے دسویں مسئلے بارے اک مختصر دستاویزی فلم جارج سِسری نے ہدایت کیتی \citep{Csicsery2016}۔ رابنسن نال یوری ماتیاسےوِچ دے سانجھے کم تے اوہدے کم اُتے رابنسن دے اثر بارے مختصر یادداشت لئی \citep{Matijasevich1992} ویکھو۔
\end{reading}

\end{document}
